\documentclass[11pt]{article}
\usepackage[a4paper,margin=1in]{geometry}
\usepackage[T1]{fontenc}
\usepackage[utf8]{inputenc}
\usepackage{lmodern,microtype}
\usepackage{amsmath,amssymb,amsthm,mathtools}
\usepackage{enumitem,xcolor,hyperref}
\usepackage[nameinlink,capitalise,noabbrev]{cleveref}
\hypersetup{colorlinks=true,linkcolor=blue!60!black,citecolor=blue!60!black,urlcolor=blue!60!black}

\newtheorem{definition}{Definition}[section]
\newtheorem{theorem}[definition]{Theorem}
\newtheorem{proposition}[definition]{Proposition}
\newtheorem{corollary}[definition]{Corollary}
\newtheorem{conjecture}[definition]{Conjecture}
\newtheorem{remark}[definition]{Remark}

\newcommand{\C}{\mathbb C}
\newcommand{\Chat}{\widehat{\mathbb C}}
\newcommand{\Oterm}{\mathcal O}
\DeclareMathOperator{\Crit}{Crit}
\DeclareMathOperator{\CV}{CV}

\title{Equal-Value Pole Monodromy for Anchored Pandrosion\\[3pt]
\large Branched coverings, moving pole sheets, and basin bifurcation conjectures}
\author{Ivan Besevic\\Independent Mathematical Research}
\date{May 2026}

\begin{document}
\maketitle

\begin{abstract}
Anchored Pandrosion maps have poles at equal-value points \(F(z)=F(a)\), not at Newton critical points. This makes the anchor a global dynamical parameter: moving \(a\) moves the pole skeleton through the level-set covering of \(F\). This paper formulates the covering theory. For a polynomial \(F:\Chat\to\Chat\) of degree \(d\), let
\[
        \mathcal P=\{(a,z):F(z)=F(a),\ z\ne a\}.
\]
Over the regular anchor set \(A^\circ=\{a:F(a)\notin \CV(F)\}\), the projection \(\mathcal P\to A^\circ\) is a \((d-1)\)-sheeted covering. Its monodromy is the ordinary inverse-branch monodromy of \(F\), with the distinguished anchor branch removed. Near a simple critical value, two equal-value sheets collide with square-root ramification, producing the natural local model for pole bifurcation. The paper states basin conjectures linking changes in anchored Pandrosion basins to this pole monodromy together with the critical orbits of the associated rational map.
\end{abstract}

\paragraph{Scope and status.}
The covering statements are elementary consequences of one-variable complex analysis. The basin statements are conjectural and belong to the complex-dynamical side of the Pandrosion programme.

\section{Equal-value pole sets}

Let \(F\) be a nonconstant polynomial of degree \(d\ge2\), viewed as a holomorphic map \(F:\Chat\to\Chat\). For a finite anchor \(a\) with \(F(a)\ne0\), the anchored equal-value pole map is
\[
        P_{F,a}(z)=a-\frac{F(a)(a-z)}{F(a)-F(z)}.
\]
Its non-removable poles lie in
\[
        \{z\ne a:F(z)=F(a)\}.
\]

\begin{definition}[Equal-value pole correspondence]
The equal-value pole correspondence is
\[
        \mathcal P_F=\{(a,z)\in\Chat^2:F(z)=F(a),\ z\ne a\}.
\]
The projection to the anchor coordinate is denoted
\[
        \pi_a:\mathcal P_F\to\Chat,\qquad (a,z)\mapsto a.
\]
\end{definition}

Let
\[
        \CV(F)=F(\Crit(F))
\]
be the finite set of critical values, including the value at infinity when appropriate. Define the regular anchor set
\[
        A^\circ=\{a\in\Chat:F(a)\notin\CV(F)\}.
\]

\section{Covering theorem}

\begin{theorem}[Equal-value pole covering]\label{thm:cover}
Over \(A^\circ\), the map
\[
        \pi_a:\mathcal P_F|_{A^\circ}\to A^\circ
\]
is a \((d-1)\)-sheeted unramified covering.
\end{theorem}

\begin{proof}
If \(a\in A^\circ\), then \(v=F(a)\) is a regular value of \(F\). Hence \(F^{-1}(v)\) consists of \(d\) distinct points. One of them is \(a\), and the remaining \(d-1\) points are the non-anchor equal-value poles.

Locally, each preimage of a regular value is a holomorphic inverse branch of \(F\). Since the anchor branch is also nonsingular and separated from the other branches, deleting the diagonal branch \(z=a\) leaves \(d-1\) holomorphic local sheets. Thus \(\pi_a\) is locally a disjoint union of biholomorphic sheets.
\end{proof}

\begin{definition}[Equal-value pole monodromy]
Choose \(a_0\in A^\circ\). Analytic continuation of the \(d-1\) non-anchor pole branches along loops in \(A^\circ\) gives a representation
\[
        \mu_F:\pi_1(A^\circ,a_0)\to S_{d-1},
\]
called the equal-value pole monodromy.
\end{definition}

\begin{proposition}[Relation with inverse-branch monodromy]\label{prop:ordinary}
The equal-value pole monodromy is the ordinary monodromy of the covering \(F:\Chat\setminus F^{-1}(\CV(F))\to\Chat\setminus\CV(F)\), pulled back along the anchor path \(a\mapsto F(a)\), with the analytically continued anchor branch distinguished and removed.
\end{proposition}

\begin{proof}
Along a loop \(a(t)\subset A^\circ\), the values \(v(t)=F(a(t))\) avoid \(\CV(F)\). The full inverse image \(F^{-1}(v(t))\) moves by ordinary inverse-branch monodromy. One branch is \(a(t)\) itself. The remaining branches are precisely the non-anchor equal-value poles.
\end{proof}

\section{Ramification and pole collisions}

\begin{theorem}[Simple critical local model]\label{thm:ramification}
Let \(c\) be a simple critical point of \(F\), and put \(v_0=F(c)\). In a local coordinate \(\xi=z-c\),
\[
        F(z)=v_0+\alpha \xi^2+\Oterm(\xi^3),\qquad \alpha\ne0.
\]
Therefore, for anchors with \(F(a)=v\) near \(v_0\), two equal-value sheets have the local form
\[
        z_\pm(v)=c\pm \sqrt{\frac{v-v_0}{\alpha}}+\Oterm(v-v_0).
\]
Analytic continuation around \(v_0\) transposes these two sheets.
\end{theorem}

\begin{proof}
The expansion follows from \(F'(c)=0\) and \(F''(c)\ne0\). The Weierstrass preparation theorem, or direct inversion of the quadratic leading term, gives the two square-root branches. A loop around \(v_0\) changes the sign of the square root, exchanging the branches.
\end{proof}

\begin{corollary}[Pole-root collision criterion]\label{cor:collision}
Equal-value pole sheets can collide only above critical values of \(F\), or with the anchor branch when the diagonal branch becomes nonregular. Thus critical values are the ramification skeleton for moving Pandrosion poles.
\end{corollary}

\section{Model example}

\begin{proposition}[Monomial monodromy]\label{prop:monomial}
For \(F(z)=z^d-1\), every finite nonzero anchor \(a\) has equal-value poles
\[
        z_j=\zeta^j a,\qquad j=1,\ldots,d-1,
\]
where \(\zeta=e^{2\pi i/d}\). A loop of \(a\) around \(0\) cyclically rotates these sheets.
\end{proposition}

\begin{proof}
The equal-value equation is \(z^d-1=a^d-1\), hence \(z^d=a^d\). The solutions are \(z=\zeta^j a\). Continuing \(a\) around \(0\) continues this labeled root system around the cyclic covering.
\end{proof}

\begin{remark}
This recovers the cyclotomic pole geometry of the monomial Pandrosion paper as the simplest monodromy model.
\end{remark}

\section{Dynamical consequences}

The anchored Pandrosion map is a rational map
\[
        P_{F,a}(z)=\frac{F(a)z-aF(z)}{F(a)-F(z)}
\]
after removable factors are cancelled. Its poles move by \(\mathcal P_F\), while its critical points satisfy the secant-tangent equation
\[
        F(a)-F(z)-(a-z)F'(z)=0.
\]
Thus basin changes can occur when pole sheets ramify, when critical orbits cross moving pole barriers, or when root multipliers cross the unit circle.

\begin{conjecture}[Pole monodromy basin bifurcation conjecture]
For a polynomial \(F\) with simple critical values, topological changes in the anchored Pandrosion basins occur only when the anchor path crosses the discriminant set determined by equal-value pole ramification, critical-orbit collision with moving poles, or loss of root attraction.
\end{conjecture}

\begin{conjecture}[Monodromy signature of basin cells]
For generic polynomial families, loops in anchor space with nontrivial equal-value pole monodromy induce nontrivial permutations of basin boundary components of the anchored Pandrosion rational map.
\end{conjecture}

\section{Conclusion}

Equal-value poles are not incidental singularities. They form a natural covering over anchor space, and their monodromy is inherited from the level-set covering of \(F\). This makes anchored Pandrosion a dynamical probe of the ramification geometry of \(F\): moving the anchor moves the pole skeleton, and the basin geometry should record that motion.

\begin{thebibliography}{9}
\bibitem{equalvalue}
I. Besevic, \emph{Complex Dynamics of Equal-Value Pole Maps}, manuscript, May 2026.
\bibitem{general}
I. Besevic, \emph{Anchored Pandrosion for General Analytic Equations}, manuscript, May 2026.
\bibitem{milnor}
J. Milnor, \emph{Dynamics in One Complex Variable}, Princeton University Press, 2006.
\bibitem{beardon}
A. F. Beardon, \emph{Iteration of Rational Functions}, Springer, 1991.
\end{thebibliography}

\end{document}
